%===========================================================
%  MEDISEG — 24-Hour Grand Finale Strategy & Idea-PPT Review
%  Slide-by-slide (1–6) with expected questions & answers
%
%  Compile:  pdflatex Hackathon24hrPlan.tex
%===========================================================
\documentclass[11pt, a4paper]{article}

%------------- Packages --------------
\usepackage[margin=2.2cm]{geometry}
\usepackage[utf8]{inputenc}
\usepackage[T1]{fontenc}
\usepackage{amsmath, amssymb}
\usepackage{xcolor}
\usepackage[colorlinks=true, linkcolor=teal!75!black, urlcolor=teal!75!black]{hyperref}
\usepackage{enumitem}
\usepackage{booktabs}
\usepackage{tabularx}
\usepackage{longtable}
\usepackage{fancyhdr}
\usepackage[most]{tcolorbox}

%------------- Colours --------------
\definecolor{mediseg}{HTML}{14b8a6}
\definecolor{medisegdark}{HTML}{0f766e}
\definecolor{slate}{HTML}{334155}

%------------- Page style --------------
\pagestyle{fancy}
\fancyhf{}
\fancyhead[L]{\small\textcolor{mediseg}{\textbf{MediSeg -- 24h Strategy \& PPT Review}}}
\fancyhead[R]{\small\textcolor{slate}{SIH 2026 \textperiodcentered\ Autodesk}}
\fancyfoot[C]{\small\thepage}
\renewcommand{\headrulewidth}{0.4pt}
\headrulecolor{mediseg!60!white}

%------------- Boxes --------------
\newtcolorbox{qabox}{
  colback=teal!6!white, colframe=mediseg, arc=3pt, boxrule=0.8pt
}
\newtcolorbox{flagbox}{
  colback=amber!12!white, colframe=orange!85!black,
  title=\textbf{Action Needed}, fonttitle=\bfseries, arc=3pt, boxrule=0.8pt
}
\newtcolorbox{warnbox}{
  colback=red!6!white, colframe=red!70!black,
  title=\textbf{Rule Reminder}, fonttitle=\bfseries, arc=3pt, boxrule=0.8pt
}

\setlength{\parskip}{3pt}

\begin{document}

%------------- Title --------------
\begin{titlepage}
\centering
\vspace*{2.4cm}
{\color{teal!50!white}\rule{\textwidth}{1pt}}\\[0.8em]
{\Huge\bfseries\textcolor{mediseg}{MediSeg}}\\[0.5em]
{\LARGE 24-Hour Grand Finale Strategy}\\[1.2em]
{\Large \& Idea-PPT Review (Slides 1--6)}\\[0.4em]
{\color{teal!50!white}\rule{\textwidth}{1pt}}\\[3em]
{\large Six-phase execution plan for the 24-hour hackathon window,\\[0.3em]
plus a slide-by-slide audit of the submitted idea deck with\\[0.3em]
expected judge questions and model answers.}\\[3.5em]

{\small Team \textbf{Maverick} \textperiodcentered\ PS ID \textbf{SIH26115}}
\vfill
{\small Version 1.0 \textperiodcentered\ September 2026}
\end{titlepage}

\tableofcontents
\clearpage

%===========================================================
\section{Reality Check: ``We have done 100\%''}

What is already \textbf{complete}:
\begin{itemize}[leftmargin=1.6em, itemsep=2pt]
  \item 6-slide idea PPT (Basic Details, Innovation, Technical Approach,
        Feasibility, Impact, References).
  \item Fully working client-side web prototype (portal, intro, login,
        interactive simulator, concept deck) deployed on Vercel.
  \item Strong-password gate, BMWM 2018 segregation simulation, traceability
        ledger, compliance report.
  \item LaTeX user guide and implementation report.
  \item Proper Git contribution setup (both repositories on the owner's GitHub
        identity).
\end{itemize}

What still belongs to the \textbf{24-hour window} (per the rules, the Fusion
product design may \emph{only} be created during the Grand Finale):
\begin{itemize}[leftmargin=1.6em, itemsep=2pt]
  \item The full Autodesk Fusion model (no pre-designed files allowed).
  \item Generative Design study, Simulation/FEA, Motion Study, Exploded
        Assembly, hi-res renders.
  \item Public link of the finished Fusion model.
  \item Final-project PPT upgrades and live demo polish.
  \item Submission package: marking-criteria table, faculty (SIH SPOC) form.
\end{itemize}

%===========================================================
\section{The 24-Hour Plan -- Six Phases}

\begin{tabularx}{\textwidth}{llX}
  \toprule
  \textbf{Phase} & \textbf{Timeline} & \textbf{Objective / Key actions}\\
  \midrule
  1 & Hours 0--1 & \textbf{Environment \& Orientation.} Confirm Fusion access for
    all members, open the Vercel demo and code repos, assign roles (CAD lead,
    simulation lead, presentation lead, demo operator). Snapshot of the
    \emph{done} deliverables so nothing is re-worked.\\
  \midrule
  2 & Hours 1--5 & \textbf{Fusion Core Design.} Parametric design of the rover:
    chassis, body, 5-compartment drum, rotating dispenser, drive module.
    Keep everything inside Fusion cloud only. Update design timeline in the
    model parameters.\\
  \midrule
  3 & Hours 5--8 & \textbf{Simulation + Generative Design.} FEA load study on
    chassis and drum rack; generative design study for structural
    optimization; capture results snapshots.\\
  \midrule
  4 & Hours 8--12 & \textbf{Motion, Assembly, Renders.} Joints + motion study
    (patrol/dumping animation), exploded assembly view, high-resolution
    renders (studio lighting).\\
  \midrule
  5 & Hours 12--16 & \textbf{Final Project PPT + Demo.} Upgrade idea deck into
    the Grand-Finale PPT: problem, solution, Fusion evidence, business model;
    screen-record the live web prototype; generate the public Fusion share
    link; insert marking-criteria table.\\
  \midrule
  6 & Hours 16--20/24 & \textbf{Rehearsal, Q\&A Drills, Submission.} Full
    run-through (6 min talk + live demo), judge Q\&A practice using the
    questions in Part~II, verify public link and uploads, submit; keep the
    remaining hours as a buffer for Vercel/link/app fixes.\\
  \bottomrule
\end{tabularx}

\begin{flagbox}
Time-box each phase. Finish with at least 2--4 buffer hours. Pre-load Fusion
tutorials and keyboard shortcuts to avoid mid-hackathon learning pauses.
\end{flagbox}

\begin{warnbox}
\textbf{Rules that never change during the 24 hours:}
\begin{itemize}[leftmargin=1.6em, itemsep=1pt]
  \item No pre-designed design files -- the Fusion model must be created only
        inside the Grand Finale window using Autodesk Fusion.
  \item The submitted design must not be copied/taken from another source.
  \item AI-generated content is not allowed in the submitted work -- the PPT,
        sketches, and Fusion model must be authored by the team.
\end{itemize}
\end{warnbox}

%===========================================================
\section{Idea-PPT Review -- Slide by Slide (1--6)}

For each slide: what it contains, a quick quality check, and
\textbf{expected judge questions with model answers}.

\subsection{Slide 1 -- Basic Details}
\textbf{On the slide:} PS ID SIH26115, PS Title ``Design and Develop a Smart
Mobile Medical-Waste Collection and Segregation System'', Theme
MedTech / BioTech / HealthTech, PS Category Software, Team \emph{Maverick}.
\begin{flagbox}
Add the \textbf{Team ID} and the names/leader of all members; double-check the
PS ID and category match the official listing.
\end{flagbox}

\begin{qabox}
\textbf{Q1.} What is your problem-statement ID and category?\\
\textbf{A1.} SIH26115, Software category -- a simulation-first solution built
with a modern web stack, supported by an Autodesk Fusion hardware design.

\textbf{Q2.} The problem mentions an \emph{Autodesk Fusion} product lifecycle --
why is your PS category ``Software''?\\
\textbf{A2.} For the idea stage we deliver a fully working software simulation
(no design files are allowed before the Grand Finale). The physical product
design is built from scratch in Autodesk Fusion during the finale, satisfying
the full concept-to-manufacturing lifecycle.

\textbf{Q3.} Who is your team?\\
\textbf{A3.} Team Maverick --- [name each member and role, Team ID].
\end{qabox}

\subsection{Slide 2 -- Innovation and Uniqueness}
\textbf{On the slide:} 5 challenges (hazardous manual handling, improper
segregation/human error, autonomous hospital navigation, AI recognition and
correct segregation, end-to-end traceability/compliance); 5 proposed-solution
capsules (simulated LiDAR/SLAM patrol, AI vision scanner with 15 items and
YOLOv8 classification logic, virtual BMWM 2018 segregation engine, digital
traceability ledger, client-side security/access gate); 4 ``how it addresses
the problems'' points.

\textbf{Check:} strong innovation story -- autonomous rover \emph{plus}
legally-compliant segregation \emph{plus} no-infrastructure deployment.

\begin{qabox}
\textbf{Q1.} How is your solution different from existing waste pushcarts?\\
\textbf{A1.} It replaces manual handling with an autonomous, battery-electric
rover that navigates autonomously, classifies waste with AI vision, routes it
into the five BMWM 2018 colour-coded compartments, and records a digital chain
of custody -- all while being a zero-infrastructure, client-side application.

\textbf{Q2.} The AI scanner is a \emph{simulation} -- is that enough?\\
\textbf{A2.} At the idea stage we simulate YOLOv8 + CNN classification with a
confidence model to prove the logic and the UI. The production path uses an
on-device YOLO model; the Fusion product design includes the camera/vision
mount and compute to realise it.

\textbf{Q3.} Why build it fully client-side with no backend?\\
\textbf{A3.} Instant deployment on any static host, zero server cost, data
stays local (privacy), and it scales horizontally across hospital networks;
the optional IoT telemetry layer can be added later without re-building the UI.
\end{qabox}

\subsection{Slide 3 -- Technical Approach}
\textbf{On the slide:} ``Technology Stack / User Interface'' (the stack grid is
image-based -- no extractable text) and a 4-step Methodology and Process.

\begin{flagbox}
\textbf{Verify the stack grid visually.} It should list: Next.js 14 (App
Router), React 18, Tailwind CSS 3.4, static export
(\texttt{output: 'export'}), Vercel (deploy), Git/GitHub (versioning). If any
of these are missing on the slide, add them -- judges expect the stack.
\end{flagbox}

\begin{qabox}
\textbf{Q1.} Why Next.js static export instead of a server?\\
\textbf{A1.} Static export gives a fast, portable, server-free app that can be
hosted anywhere (Vercel, Netlify), keeps costs near zero, and is inherently
scalable -- ideal for a demo that must also be production-lean.

\textbf{Q2.} How does classification work at runtime?\\
\textbf{A2.} A confidence model draws classification certainty and a 2\,\%
misclassification trigger routes waste to the safe (infectious) bin as a
fail-safe; the pipeline animates detect $\rightarrow$ classify $\rightarrow$
segregate stages. In production this becomes a real YOLO model on edge
hardware.

\textbf{Q3.} What about security?\\
\textbf{A3.} A reviewable client-side gate protects the prototype: a
five-rule strong-password validator (length, upper, lower, digit, special)
and a session stored in the browser. A production release would add a
real authentication backend.
\end{qabox}

\subsection{Slide 4 -- Feasibility and Viability}
\textbf{On the slide:} the content is visual (graphics/icons) with no
extractable text.
\begin{flagbox}
\textbf{Check the graphics are legible.} Ensure the slide clearly shows the
four feasibility pillars (technical, operational, financial, regulatory).
If the current graphic is unclear, replace with a simple four-quadrant table
-- substance over decoration.
\end{flagbox}

\begin{qabox}
\textbf{Q1.} Is this feasible without real hardware?\\
\textbf{A1.} For the idea stage, yes -- the simulation is complete and
deployed. The Fusion design in the finale demonstrates manufacturability with
focus on cost, materials (ABS+steel), and serviceability (hot-swap battery).

\textbf{Q2.} Financial viability?\\
\textbf{A2.} Subscription SaaS tiered by facility beds and transaction volume,
plus enterprise licensing for regulators; low marginal cost because the
software is static-cloud and the hardware is a commodity electric platform.

\textbf{Q3.} Regulatory viability?\\
\textbf{A3.} Built directly on BMWM 2018 colour-coding and produces a
CPCB-ready digital manifest with an auditable, cradle-to-grave ledger.
\end{qabox}

\subsection{Slide 5 -- Impact and Benefits}
\textbf{On the slide:} Business Model (target customers, value-added services,
scalability, financial model) and four audience impacts (hospital
administrators, healthcare workers/sanitation staff, regulators CPCB/state
boards, engineering students/evaluators).

\textbf{Check:} strong customer framing; consider adding a headline impact
statistic (e.g., ``$\sim$98\,\% classification, zero-touch handling,
40\,\% lower infection risk target'').

\begin{qabox}
\textbf{Q1.} Who is your customer and how do you make money?\\
\textbf{A1.} Hospitals and biomedical-waste operators needing BMWM 2018
compliance; regulators needing audit trails. Revenue via tiered SaaS
subscriptions and enterprise licensing for regulatory boards.

\textbf{Q2.} What concrete impact do you claim?\\
\textbf{A2.} Reduced manual touchpoints and needle-stick risk for sanitation
staff, fewer non-compliance penalties for hospitals, transparent verifiable
audits for regulators, and a reusable reference architecture for builders of
simulation tools.

\textbf{Q3.} How do you prove the impact metrics?\\
\textbf{A3.} The live dashboard reports accuracy, capacity, and compliance
checks in real time; the ledger gives measurable audit evidence. Pilot
evidence targets would follow in a hospital trial.
\end{qabox}

\subsection{Slide 6 -- Research and References}
\textbf{On the slide:} a sample UI image developed by the team, and three
references (PRISMA systematic review on real-time tracking technologies in
biomedical waste management; AI-powered biomedical segregation and recycling
with IoT; IoT-enabled biomedical smart waste bin monitoring using YOLO).

\textbf{Check:} confirm the three links are accessible and ready for judges;
keep the sample-UI image crisp.

\begin{qabox}
\textbf{Q1.} What research informed your design?\\
\textbf{A1.} The PRISMA review guided our tracking/traceability architecture;
the IoT-YOLO bin-monitoring paper validated vision-based waste
classification; the AI-segregation study informed the compartment routing
logic -- combined with BMWM 2018 regulatory requirements.

\textbf{Q2.} Which reference did you learn the most from and why?\\
\textbf{A2.} The IoT-enabled YOLO monitoring study, because it confirmed that
real-time computer vision can drive accurate waste categorisation at the
source -- directly supporting our camera-scan pipeline.

\textbf{Q3.} Is there a competitive analysis?\\
\textbf{A3.} Our positioning: existing solutions cover monitoring or bins
separately; MediSeg unifies autonomous collection, AI classification,
segregation, and audit traceability in a single no-infrastructure platform.
\end{qabox}

%===========================================================
\section{Marking Criteria -- Proposed Framework}

The statement requires a marking-criteria table to be attached. Confirm
against the official SIH 2026 rubric, then insert into the idea form/PPT.

\begin{tabularx}{\textwidth}{lcc}
  \toprule
  \textbf{Criterion} & \textbf{Weight} & \textbf{Covered in our deck}\\
  \midrule
  Problem understanding & 15 & Slide 2 --- challenges\\
  Innovation \& uniqueness & 20 & Slide 2 --- solution capsules\\
  Technical approach & 20 & Slide 3 --- stack + methodology\\
  Feasibility \& viability & 15 & Slide 4 --- feasibility pillars\\
  Impact \& benefits & 15 & Slide 5 --- business model + audience\\
  Presentation \& communication & 10 & Slides 1, 6 + live demo\\
  Regulatory alignment & 5 & BMWM 2018, CPCB manifest\\
  \midrule
  \textbf{Total} & \textbf{100} & \\
  \bottomrule
\end{tabularx}

\begin{flagbox}
Verify the weights against the official SIH 2026 marking criteria before
attaching; this is a proposed framework only.
\end{flagbox}

%===========================================================
\section{Submission Checklist}

\begin{itemize}[leftmargin=1.6em, itemsep=2pt]
  \item 6--7 slide idea PPT (complete) with conceptual sketches, research,
        and images.
  \item Attached marking-criteria table.
  \item Faculty (SIH SPOC) mandatory form requested and confirmed.
  \item Grand Finale (24h): Fusion model created fresh inside the event,
        public share link generated, generative design + simulation + motion +
        exploded view + hi-res renders captured.
  \item Final-project PPT and recorded demo exported; live link verified
        (\url{https://mediseg-sih2026.vercel.app}).
  \item No pre-designed files; no AI-generated content in submitted work.
\end{itemize}

\end{document}